\begin{table}[htbp]
\centering
\caption{E2算法消融结果}
\label{tab:e2_ablation}
\small
\begin{tabular*}{\linewidth}{@{\extracolsep{\fill}}lrrrrrr@{}}
\toprule
算法臂 & Best（元） & Avg（元） & Gap（\%） & 车辆数 & CPU（min） & MV胜/平/负 \\
\midrule
O（纯距离开源） & 3609.597 & 3615.557 & 1.919 & 12.111 & 0.759 & 354/46/5 \\
HGS-F & 3587.524 & 3592.472 & 1.059 & 12.165 & 1.852 & 307/98/0 \\
HGS-E & 3558.578 & 3563.655 & 0.069 & 12.207 & 2.081 & 125/280/0 \\
HGS-M & 3558.450 & 3566.005 & 0.085 & 12.230 & 2.104 & 114/291/0 \\
\textbf{MV-HGS-SP} & \textbf{3556.687} & \textbf{3559.608} & 0.000 & 12.178 & 6.209 & 0/405/0 \\
\bottomrule
\end{tabular*}
\tabnote{Best为81个算例各自5个种子最优成本的等权均值，Avg为405个实例--种子单元成本的等权均值。Gap$=100\times\mathrm{mean}[(C_{\rm arm}-C_{\rm MV})/C_{\rm arm}]$，正值表示MV-HGS-SP成本较低；车辆数为每个封存方案中唯一vehicle\_id的数量。CPU仅作描述：O臂与其余四臂非同批、非同机、非同停止规则且非等算力。}
\end{table}
